\documentclass[instructions.tex]{subfiles}
\begin{document}
 
\chapter{Il est dangereux pour le salut d’être riche.}
 
Nous allons montrer les dangers des richesses, et l’usage qu’on doit faire des richesses et de la pauvreté.
 
\primo Oh! qu’il est difficile à ceux qui aiment l’argent d’entrer dans le ciel! dit Jésus-Christ. Il est plus facile de faire passer un chameau par le trou d’une aiguille, qu’un riche par la porte du ciel. Pourquoi? parce qu’il est difficile d’être riche sans aimer les richesses. Or, dès qu’on les aime et qu’on s’y attache, elles font oublier ce qu’on doit à Dieu, ce qu’on doit à soi-même, ce qu’on doit aux autres.
 
Ce qu’on doit à Dieu, ce qu’on doit d’abord à Dieu. Est-on riche, dit le prophète Osée, on a trouvé son idole et son Dieu\footnote{Diyes effectus sunt, inveni idolum mini, Os. 12.}. Il y a cette différence entre l’idolatre et le riche, que l’idolatre adore l’or en statue, et que le riche l’adore en monnoie et en fonds. En effet, loin de sacrifier à Dieu son temps, il le sacrifie à tous ses biens, ne pense qu’à ses revenus et à ses pertes. Oubliant que Dieu est l’auteur de tout bien, il n’a point de confiance en sa providence, et craint toujours de manquer. De là les emportemens quand il perd ou qu’on lui fait tort; les menaces et les chagrins quand les saisies sont fâcheuses.
 
Saint Paul dit que la cupidité, l’amour des richesses a fait perdre la foi à plusieurs\footnote{Erraverunt à fide. S. Paul.}. C’est en vain qu’un pasteur annonce au riche les vérités de la foi, et l’obligation de donner de bons exemples. Le son d’un écu, dit saint Ambroise, fait bien plus d’impression sur son cœur que le son de la parole de Dieu. Pourquoi les païens se moquaient-ils de la doctrine de Jésus-Christ? C’est parce qu’ils étaient avares, attachés aux richesses.
 
On se croit dispensé d’obéir aux lois de Dieu et de l’Église, aussitôt qu’on a du bien. La religion n’a point ordinairement de plus grands ennemis que certains riches. Tout ce qu’on peut perdre pour l’amour de Dieu, ce qu’on retranche pour la décoration du lieu saint, pour les établissemens de piété, pour l’instruction des âmes, leur semble perdu. Pourquoi? disent-ils avec l’avare Judas\footnote{ Utquid perditio hæc? Matth. 26. 8.}? Il n’est donc rien de plus inique que d’aimer l’argent, dit l’Ecclésiastique \footnote{ Nihil est iniquiùs quam amare pecuniam. Eccli. 10.}.
 
2. Un homme attaché à un objet, oublie tout le reste. Un riche attaché à ses biens, oublie tout, jusqu’à oublier ce qu’il doit à lui-même et à sa famille. Ses revenus, ses fonds, ses animaux, lui tiennent bien plus à cœur que son salut, que l’éducation de ses enfans, de ses domestiques. S’il donne quelque attention à sa famille, il a bien plus d’empressement à l’enrichir qu’à la sanctifier. Deux grands abus se trouvent ordinairement dans les richesses: c’est l’acquisition des biens, et l’emploi qu’on en fait.
 
Les chemins qu’on prend pour arriver à la fortune sont comme ces chemins écartés où l’on met un signal pour avertir les passans: Ici un homme a été volé, Ici un autre a été blessé, un autre assassiné; prenez garde à vous.
 
Si ceux qui veulent acquérir du bien écoutaient le signal de la conscience, elle leur dirait: Prends garde à toi. La route que tu tiens en a fait damner plusieurs: les uns, par la rapine, par l’usurpation des fonds, l’anticipation sur les voisins, la dégradation des forêts; les autres, par les procès, par la chicane et la mauvaise foi; d’autres, par l’usure, par les concussions, par l’exaction de certains droits, etc. Mais on veut être riche, la cupidité étouffe la voix de la conscience.
 
Ce n’est pas tout: quel usage fait-on des grandes richesses? N’en voit-on pas qui changent même d’y toucher, qui refusent à leurs femmes le nécessaire, à leurs enfans l’établissement, aux domestiques le salaire, à eux-mêmes la subsistance, et se rendent un objet d’exécration, à Dieu et à tout le voisinage? D’autres riches sont des hommes et à sans conscience qui emploient sans honte et des moyens pour acquérir leurs richesses et leur crédit, parce qu’ils ont la crainte de Dieu; mais ceux qui n’ont pas cette crainte du Seigneur, semblent n’être riches que pour réunir en eux tous les vices, la vanité par des dépenses fastueuses en luxe, en jeux, en repas somptueux, grand train, grand nombre de domestiques vicieux et fainéans, 2. L’ambition: jamais contens de leur état, ils se servent des moyens les plus odieux, et de leur crédit, pour supplanter, pour s’élever. 3. La volupté; leur corps, nourri dans la mollesse, est un cloaque d’impureté. Combien de personnes du sexe séduites par l’argent, par les présens, par la violence! leur maison, sans règle, est peut-être la source et la pépinière des désordres d’une paroisse et du voisinage. 4. L’impiété, par les discours qu’ils tiennent contre la religion, par le mépris qu’ils font des sacremens et des mystères de Dieu, par les livres pernicieux qu’ils se procurent et qu’ils répandent, par les scandales qu’ils souffrent et qu’ils autorisent, etc.
 
Ils se persuadent que tout leur est permis, parce qu’on les flatte et qu’on n’ose les contredire; ils se croient sans reproche, parce que personne, pas même les pasteurs de l’Église, n’osent les reprendre. Oh! dans quels abîmes les richesses n’entraînent-elles pas ceux qui n’en usent pas selon Dieu!
 
3. Si les richesses font oublier ce qu’on doit à soi-même et à sa conscience, elles font encore plus oublier ce qu’on doit à autrui. Un riche puissant, qui n’a pas la crainte de Dieu, ne pense qu’à dominer; il se persuade que les autres ne sont que pour lui. Plein d’orgueil, il veut s’égaler à ceux qui lui sont supérieurs en naissance; c’est un insensé à qui la fortune a tourné la tête, qui ne se souvient plus de ce qu’il est. Rempli de lui-même, il n’a que de l’indifférence pour ses égaux et du mépris pour ceux qu’il croit au-dessous de lui; il ne prend pas garde qu’il méprise des gens qui, quoique moins riches, valent mieux que lui. On dirait qu’il n’est riche que pour se rendre lui-même plus méprisable et se faire haïr.
 
Le pauvre et le créancier lui sont un objet d’horreur; il n’a pour le misérable et le débiteur que des entrailles de fer. Par qui les pauvres sont-ils rebutés et surchargés? n’est-ce pas ordinairement par les riches? Qu’est-ce qui fait la fortune de certains riches? l’usurpation des biens du pauvre. De quoi sont bâties leurs maisons superbes? de quoi sont-elles meublées? Des sueurs et du sang des pauvres, dit Jérémie\footnote{In atriis tuis inventus est sanguis animarum pauperum, Jer. 2.}. Doit-on s’étonner si Jésus-Christ a donné sa malédiction aux riches, et si l’Écriture parle si hautement des malheurs qui les attendent à leur dernière heure? Le riche mourut, dit l’Évangile, et il fut enseveli dans l’enfer. Voilà le panégyrique que le Saint-Esprit fait d’un riche après sa mort.
 
\secundo Dieu ne condamne cependant pas les richesses, mais le cœur qui s’y attache, dit saint Augustin\footnote{Nod divitias damnat, sed cor appositum damnat.}. Salomon a été riche, et pour cela il n’était pas coupable. Abraham, Job, saint Louis, ont été riches, mais ils se sont sanctifiés dans leurs richesses.
 
Un riche qui veut se sauver, doit, selon l’avis de l’apôtre saint Jacques, regarder ses richesses comme des sujets de larmes, et pousser des gémissemens sur les malheurs dont il est menacé\footnote{Agite nunc, divites; plorate, ululantes in miseriis vestris quæ advenient vobis. S. Jac.}. Il devrait rougir d’avoir si peu de conformité avec Jésus-Christ, qui, quoique maître du monde, n’avait pas même où reposer sa tête.
 
Tout ce qui peut consoler un riche, c’est qu’il peut protéger le pauvre, l’assister, le secourir, et racheter ses péchés par d’abondantes aumônes. Pour en venir à la pratique, contentez-vous des richesses que le Seigneur vous envoie par la succession de vos pères ou par un travail légitime; possédez-les sans attache; servez-vous-en avec modération; répandez-les avec libéralité; perdez-les avec générosité et sans chagrin; et d’un danger de damnation vous en ferez un moyen de salut.
 
\tertio La pauvreté n’a pas le danger des richesses; elle est un moyen de salut quand on la souffre avec patience et qu’on met sa confiance en Dieu. Ne craignez rien, mon fils, disait Tobie; nous sommes pauvres, mais nous aurons toujours assez de biens, et nous serons toujours contens, si nous avons la crainte de Dieu.
 
La pauvreté ne laisse pas d’avoir ses écueils, elle devient un danger de damnation à celui qui n’en profite pas. Les pauvres ordinairement ne s’en occupent et ne parlent que de leurs nécessités; leurs confessions ne sont souvent qu’un récit ennuyeux et affecté de leurs misères; ils sont beaucoup plus en peine des maux du corps que de ceux de l’âme, et pensent bien plus à se tirer de l’indigence qu’à se tirer de l’enfer. Impatiens, avidés, murmuratcurs; ingrats, envieux, voleurs, fourbes, traîtres, menteurs, paresseux, ivrognes: voilà les désordres dans lesquels la pauvreté entraîne un homme sans vertu. O pauvres! que vous êtes aveugles de perdre la confiance en Dieu, qui voit vos peines! Que vous êtes insensés, d’un moyen de salut, d’en faire un danger de damnation!
 
\section{Résolutions}
 
\primo Si j’ai du bien, je donnerai mon superflu aux pauvres.

\secundo Si je suis pauvre, je supporterai cette position avec patience, résignation, sans envie et sans murmure. 

\tertio Je regarderai les richesses comme des moyens de faire le bien, et comme pleines de dangers pour le salut. 

\prayer{Mon Dieu, aidez-moi à faire un saint usage de la situation où votre Providence m’a placé par rapport aux biens passagers de ce monde.}
 
\end{document}
